% Pillar-5 (MIN-REPORT-RL) related work
\section{Related Work}
\label{sec:p5-related}

\paragraph{Companion position statement.} A condensed community-facing
statement of the same eight-item standard exists as the ZVF-program
position draft \emph{``MIN-REPORT-RL: Reporting the Stack, Not the
Label''} (\texttt{zvf-program/position/min\_report\_rl.tex}). This paper
is canonical for the evidence corpus and the standard's justification;
the position draft is canonical for the community-adoption framing. The
item lists were aligned at eight on 2026-07-11.

\paragraph{Reproducibility crises in deep RL.} Our position is a deliberate
echo of \citet{henderson2018deep}, who showed that deep-RL results depend on
codebase, hyperparameters, and even the random seeds chosen, to the point
that the same algorithm from two repositories behaves as two algorithms.
\citet{engstrom2020implementation} showed that unlabelled ``code-level
optimizations'' account for much of PPO's reported advantage over TRPO---the
direct ancestor of our label-versus-stack argument.
Follow-on work built the statistical repair kit: rigorous comparison
protocols \citep{colas2019hitchhiker} and small-sample-aware aggregate
metrics \citep{agarwal2021deep}, and position work has questioned what RL
benchmarking can establish at all \citep{jordan2024benchmarking}. RL-for-LLMs
adds a new failure surface those works did not face: the sampling
distribution is produced by a serving stack complex enough to be a research
artifact in its own right, and often closed.

\paragraph{Reporting standards in NLP and ML.} \citet{dodge2019show} argued
that results are functions of compute budget and hyperparameter search, and
proposed reporting expected validation performance as a function of budget;
our manifest is the RL-for-LLM analogue, replacing ``budget'' with ``stack.''
The NeurIPS reproducibility program \citep{pineau2020improving} demonstrated
that checklists plus process shift community behavior---and also that
self-attested checklists have limits, which motivates our artifact-computed
badge (\secref{sec:p5-toolchain}). Statistical treatments of evaluation
noise \citep{miller2024erroreval, zhang2024benchvariance} quantify the
uncertainty floor against which our 12-cell head-to-head
(\secref{sec:p5-exhibit-h2h}), which is \emph{not} shown to be within noise,
should be read.

\paragraph{Evaluation-stack sensitivity.} The lm-evaluation-harness effort
documented, from years of maintenance experience, that prompt formatting,
answer extraction, and implementation details shift LLM benchmark scores by
ranking-relevant margins, and argued for shared, versioned evaluation
infrastructure \citep{biderman2024lessons}. Independent audits reach the same
conclusion for reasoning evaluations specifically: \citet{hochlehnert2025sober}
find math-reasoning results sensitive to decoding, seeds, prompt format, and
hardware/software configuration, and that many reported RL gains shrink under a
standardized protocol, and contamination studies show the data side of
the same coin \citep{zhang2024gsm1k, riddell2024contamination}. Our claim is
that RL \emph{training} inherits all of this---the verifier is an evaluator
in the loop---plus the V1--V3 vectors unique to on-policy optimization.

\paragraph{Silent framework defects and controlled comparisons.}
\citet{limozin2026sftrl} trace published gains of mixed-policy methods over
SFT-then-RL to two silent framework bugs (a DeepSpeed CPU-offload optimizer bug
affecting TRL, OpenRLHF, and LLaMA-Factory, and an OpenRLHF loss-aggregation
bug)---a stack-provenance failure no algorithm label would reveal.
\citet{liu2025tricks} re-run common RL-for-reasoning tricks in one framework
and find that conflicting published conclusions trace to differing
settings, data, and initialization. Controlled multi-algorithm studies
\citep{li2026oxrl} report scale-dependent ranking inversions on identical
infrastructure, and critic-based methods remain competitive when implemented
carefully \citep{kazemnejad2025vineppo, liu2025asyppo}, so a ``GRPO vs.\ PPO''
gap measured on one stack is a property of that stack. MIN-REPORT-RL is a
reporting-side complement: it does not prevent such defects but records the
fields needed to localize them.

\paragraph{The GRPO-variant explosion.} The label space we target is growing
quickly: GRPO \citep{shao2024deepseekmath, guo2025deepseekr1}, DAPO
\citep{yu2025dapo}, Dr.GRPO \citep{tong2025drgrpo}, GSPO \citep{qwen2025gspo},
and a long tail of descendants \citep{nan2025ngrpo,
zhang2025scaffgrpo, lu2025grpolead, bytedance2025vapo}. Analyses that unify
the family---GRPO's preference-learning reading \citep{rafailov2024direct,
wu2025grpodpo} and zero-variance-prompt recovery methods \citep{le2025rlzvp}
---support our premise that members differ by a small set of loss-form
coordinates, which is exactly what item 1 of the manifest records. Empirical
scaling studies of GRPO knobs \citep{tan2025scalingrl} likewise presuppose
the per-knob disclosure the manifest mandates. The companion \ours{} pillar
papers (P1 scaling, P2 ZVF, P3 group size, P4 length bias) provide the
measurement substrate this position paper argues should become the norm.

\paragraph{Documentation standards.} MIN-REPORT-RL follows the documentation
tradition of \emph{Model Cards} \citep{mitchell2019modelcards}, which record
training procedure, intended use, and per-slice evaluation for trained models;
\emph{Datasheets} \citep{gebru2021datasheets} and \emph{Data Statements}
\citep{bender2018datastatements}, which record dataset composition,
collection, and curation; \emph{Data Cards} \citep{pushkarna2022datacards},
which record dataset provenance and sampling for practitioners; and
FactSheets (Arnold et al., 2019, \emph{IBM J.\ Res.\ Dev.}), which frame such
documents as a supplier's declaration of conformity. These standards document
artifacts; MIN-REPORT-RL documents the training \emph{comparison}, and its
items are chosen for whether they can move or flip a reported ranking
(\secref{sec:p5-stack}) rather than for documentation completeness.
